\scanpage{20}
% archive sheet number: 19

A) Cas $n$ pair $\geq 2$, donc $C$ de rang $n$ pair aussi donc $V$ de rang impair, donc
\[
\operatorname O(f_V)\supset\operatorname{SO}(f_V)\times\mu_2,\quad\text{et}
\]
\diagram{\begin{tikzcd}[ampersand replacement=\&]
\operatorname{SO}^{\unsure{\circ}}(f_V)\arrow[r,hook]\&\operatorname{Aut}(C,Q_C)
\end{tikzcd}}

En fait, l'hom
\diagram{\begin{tikzcd}[ampersand replacement=\&]
G\arrow[r,"{\varphi_G^C}"]\&\operatorname{Aut}(C,Q_C)
\end{tikzcd}}

se factorise à travers $\operatorname{SO}(f_V)$
\diagram{\begin{tikzcd}[ampersand replacement=\&]
G\arrow[r,"{\psi_G^V}"]\&\operatorname{SO}(f_V)
\end{tikzcd}}

On notera que
\[
\psi_{\unsure{G}}^V(\underline{\tau}_i)=-\varphi_G^V(\underline{\tau}_i)=-\tau_i^V
\]

\uline{Ex} \uline{$n=2$,} \deleted{alors \unsure{$Q$}} Supposons $f_V$ lisse\deleted{\unsure{en|lisse}} i.e. $Y=Q_C$ un fibré en droites projectives i.e.
\[
\underset{\substack{\shortparallel\\\alpha_0+2}}{\beta_0},\
\underset{\substack{\shortparallel\\\alpha_1+2}}{\beta_1}
\ \text{et}\ \beta_0+\beta_1-4=\alpha_0+\alpha_1\ \text{inv.},\quad\text{alors}
\]
\[
\operatorname{SO}(f_V)\xrightarrow{\sim}\operatorname{Aut}(Q_C),
\]
donc la donnée de $C$ et $\varphi_G^C$ ou $\psi_G^V$ équivaut à \deleted{celle} du fibré en droites projectives $Y$ $(=Q_C)$, et de l'opération de $G$ dessus via $\tau_0^Y,\tau_1^Y,\tau_2^Y$. Traduisons les conditions explicitées plus haut sur $\unsure{\psi_G^C|\varphi_G^C}$ en termes de $\varphi_G^Y$ (\uline{NB} ici $\beta_0$ inv. donc $\tau_0$ \unsure{réflexion|psflexion} stricte, donc $h_0$ s'en déduit \unsure{de|du} l'\unsure{enveloppe|un} \unsure{$2$-projective|$2$-projectiv.|$2$-rejet} $C$ de $Y$)
